%%%%%%%%%%%%%%%%%%%%%%%%%%%%%%%%%%%%%%%%%%%%%%%%%%%%%%%%%%%%%%%%%%%%%%%%%%%%%
% Diffusion tech report, NeurIPS 2026 formatting (same setup as the CS336
% assignment-1/2 reports). Compile:  latexmk -pdf tech_report.tex
%
% Figures are produced by make_figures.py from results/*.json; missing ones
% render as a placeholder box so the document always compiles.
%%%%%%%%%%%%%%%%%%%%%%%%%%%%%%%%%%%%%%%%%%%%%%%%%%%%%%%%%%%%%%%%%%%%%%%%%%%%%
\documentclass{article}

\PassOptionsToPackage{numbers, compress}{natbib}
\usepackage[preprint]{neurips_2026}

\usepackage[utf8]{inputenc}
\usepackage[T1]{fontenc}
\usepackage{hyperref}
\usepackage{url}
\usepackage{booktabs}
\usepackage{amsfonts}
\usepackage{amsmath}
\usepackage{amssymb}
\usepackage{nicefrac}
\usepackage{microtype}
\usepackage{xcolor}
\usepackage{graphicx}
\usepackage{listings}

\lstset{
  basicstyle=\ttfamily\footnotesize,
  columns=fullflexible,
  keepspaces=true,
  breaklines=true,
  frame=single,
  framerule=0.4pt,
  rulecolor=\color{gray!60},
  backgroundcolor=\color{gray!8},
  xleftmargin=1.5em,
  xrightmargin=1.5em,
  aboveskip=1.2em,
  belowskip=1.2em,
}

\usepackage{caption}
\captionsetup{font=small}
\usepackage{algorithm}
\usepackage{algpseudocode}

% Figures that have not been generated yet become a visible placeholder
% instead of a hard compile error.
\newcommand{\fig}[2][width=\linewidth]{%
  \IfFileExists{#2}{\includegraphics[#1]{#2}}{%
    \fbox{\parbox[c][2.6cm][c]{0.92\linewidth}{\centering\ttfamily\footnotesize
      \detokenize{#2}\par\vspace{3pt}%
      not generated yet --- run \detokenize{python reports/make_figures.py}}}%
  }%
}

\title{From DDPM to Rectified Flow:\\
a Measured, From-Scratch Diffusion Study}

\author{%
  Yi Hou \\
  University of Chinese Academy of Sciences\\
  Beijing, China \\
  \texttt{houyi25@mails.ucas.ac.cn}
}

\begin{document}

\maketitle

\begin{abstract}
This report documents a from-scratch implementation and measurement study of the modern
diffusion stack. No \texttt{diffusers}, \texttt{flash\_attn} or pretrained generative
component is used: DDPM, DDIM, a diffusion transformer (DiT) with adaLN-Zero conditioning,
and rectified flow are all implemented in this repository, together with the training
pipeline, the sampler suite and the evaluation code. The study is organized as a controlled
2$\times$2 design (UNet vs.\ DiT $\times$ $\epsilon$-prediction vs.\ rectified flow) at two
scales, MNIST $28\times28$ and CelebA $64\times64$, with every cell trained under one frozen
protocol. Headline results: DDIM at 50 function evaluations recovers the 1000-step ancestral
FID to within 12\% on CelebA ($12.75$ vs.\ $11.37$) at $5.2\times$ the sampling throughput,
and rectified flow dominates the quality--NFE frontier on MNIST (FID $0.88$ at 64
evaluations, versus $4.32$ for the best $\epsilon$-prediction sampler at 50). The same
deterministic sampler that improves the UNet collapses on the transformer --- its FID grows
from $49$ to $672$ as NFE goes $10\rightarrow50$, against $33$ for 1000-step ancestral
sampling of the same weights --- and an $\eta$ sweep shows the failure is the missing noise
injection. On the systems side, bf16 buys $3.15\times$ step time and $1.6\times$ memory over
fp32 while \texttt{torch.compile} buys nothing ($1.00\times$), and a kernel-level profile
attributes 67\% of GPU time to matmuls and 16\% to elementwise work, which is where a fused
adaLN kernel would have to earn its keep.

\end{abstract}

%=============================================================================
\section{Introduction}
%=============================================================================

Diffusion models are usually consumed through a library: a scheduler object, a pretrained
backbone, a pipeline. This report takes the opposite route. Everything is built from
scratch --- the forward and reverse processes, the two backbones, both training objectives,
the samplers and the evaluation --- and every claim is backed by a measurement produced by
this repository.

The study asks three questions:

\begin{itemize}
  \item \textbf{RQ1 (architecture).} At a similar parameter scale, does replacing a UNet
        with a transformer (DiT) backbone help or hurt generation quality and computational
        efficiency?
  \item \textbf{RQ2 (objective and sampler).} How does a rectified-flow objective change the
        quality--NFE trade-off, and how much does sampler choice alone buy for a fixed
        trained model (DDPM ancestral vs.\ DDIM)?
  \item \textbf{RQ3 (systems).} How much training and inference throughput can be recovered
        by mixed precision, compilation and kernel-level fusion at this scale --- and where
        does a distributed strategy stop being worth it?
\end{itemize}

Section~\ref{sec:background} puts the three objectives on a common footing.
Section~\ref{sec:implementation} describes the implementation,
Section~\ref{sec:setup} the frozen protocol and the evaluation protocol,
Sections~\ref{sec:samplers}--\ref{sec:systems} the results, and
Section~\ref{sec:failures} collects the experiments that failed, which were as instructive
as the ones that worked.

\begin{table}[t]
  \caption{Results at a glance. All numbers are measured on the hardware described in
  Section~\ref{sec:setup}; the raw archives live in \texttt{results/*.json}.}
  \label{tab:glance}
  \centering
  \small
  \begin{tabular}{@{}llp{0.5\linewidth}@{}}
    \toprule
    Question & Baseline & Measured effect \\
    \midrule
    RQ1 backbone & UNet (11.7M params, 2.43 GFLOPs/sample) &
      DiT (16.5M, 1.08): worse FID at every matched budget
      (MNIST $\epsilon$: $2.79$ vs.\ $33.2$; flow: $0.88$ vs.\ $4.99$) despite $2.25\times$
      fewer FLOPs \\
    RQ2 objective & $\epsilon$-prediction &
      rectified flow: FID $0.88$ vs.\ $4.32$ at NFE 64 vs.\ 50; crossover near NFE 16 \\
    RQ2 sampler & DDPM ancestral (1000 NFE) &
      DDIM at 50 NFE: FID $12.75$ vs.\ $11.37$ on CelebA ($5.2\times$ throughput);
      collapses on the DiT ($49\!\rightarrow\!672$ as NFE $10\!\rightarrow\!50$) \\
    RQ3 systems & FP32, uncompiled &
      BF16 $3.15\times$ step time, peak memory $32.9\!\rightarrow\!21.1$\,GB, MFU
      $10.4\%\!\rightarrow\!32.7\%$; \texttt{compile} $1.00\times$ \\
    RQ3 profile & --- &
      matmul 67.1\%, elementwise 15.6\%, attention 11.1\%; the fused-adaLN target is the
      elementwise slice \\
    \bottomrule
  \end{tabular}
\end{table}

%=============================================================================
\section{Background: one path, three objectives}
\label{sec:background}
%=============================================================================

All three methods place the data and the noise at the ends of an interpolation path and
train a network to regress a target along it. They differ in the path, in the target, and in
the solver used to walk back from noise to data.

\begin{table}[t]
  \caption{The three objectives on a common footing. $\bar\alpha_t$ is the usual DDPM
  cumulative product; the flow variant uses a continuous $t\in[0,1]$.}
  \label{tab:objectives}
  \centering
  \small
  \begin{tabular}{@{}llll@{}}
    \toprule
    Method & Path $x_t$ & Target & Sampler (NFE) \\
    \midrule
    DDPM~\citep{ho2020ddpm} & $\sqrt{\bar\alpha_t}\,x_0 + \sqrt{1-\bar\alpha_t}\,\epsilon$
      & $\epsilon$ & ancestral, stochastic (1000) \\
    DDIM~\citep{song2020ddim} & same & $\epsilon$ & deterministic ODE (10--50) \\
    Rectified flow~\citep{liu2022flow,lipman2022fm} & $(1-t)\,x_0 + t\,\epsilon$
      & $v = \epsilon - x_0$ & Euler ODE (4--128) \\
    \bottomrule
  \end{tabular}
\end{table}

\begin{figure}[t]
  \centering
  \fig[width=2.65in]{figures/paths.pdf}
  \caption{The diffusion path and the rectified-flow path between data and noise. The flow
  objective regresses a constant velocity along a straight line; the diffusion objective
  traverses a curved, non-uniform-speed path.}
  \label{fig:paths}
\end{figure}

One implementation detail is easy to get wrong and worth stating: the sinusoidal timestep
embedding used by both backbones is calibrated for $t \in [0, 1000]$ (the DDPM range). The
flow code therefore keeps its path mathematics in $t\in[0,1]$ but passes $t\times1000$ to the
network; without that scaling the embedding is nearly constant and conditioning collapses.

%=============================================================================
\section{Implementation}
\label{sec:implementation}
%=============================================================================

Everything is written in PyTorch, with no generative-model library. Table~\ref{tab:models}
lists the measured sizes. Note that the parameter ranking and the FLOPs ranking disagree
across scales --- on CelebA the UNet is larger on both axes, on MNIST it has fewer
parameters but more FLOPs --- which is exactly why both are reported.

\begin{table}[t]
  \caption{Measured sizes (forward pass, one sample at the model's native resolution).
  Params and FLOPs are counted by \texttt{eval/cost.py} with
  \texttt{torch.utils.flop\_counter}, so both backbones are measured with the same tool.}
  \label{tab:models}
  \centering
  \small
  \begin{tabular}{@{}llrrrr@{}}
    \toprule
    Dataset & Backbone & Params & GFLOPs/sample & Conditioning & Attention \\
    \midrule
    MNIST $28^2$ & UNet & 11.7M & 2.43 & sinusoidal + MLP & bottleneck \\
    MNIST $28^2$ & DiT  & 16.5M & 1.05 & adaLN-Zero          & all 6 blocks (49 tokens) \\
    CelebA $64^2$ & UNet & 174.1M & 67.8 & sinusoidal + MLP & bottleneck \\
    CelebA $64^2$ & DiT  & 129.6M & 43.6 & adaLN-Zero          & all 12 blocks (256 tokens) \\
    \bottomrule
  \end{tabular}
\end{table}

\textbf{Backbones.} The UNet is a residual U-net with GroupNorm, SiLU, sinusoidal timestep
injection into every residual block, and \texttt{scaled\_dot\_product\_attention} at the
bottleneck. The DiT follows \citet{peebles2023dit}: patchify with a strided convolution,
fixed two-dimensional sincos positional embeddings, and blocks whose attention and MLP are
modulated by an adaLN-Zero MLP that is zero-initialized, so every block starts as the
identity (verified in \texttt{tests/test\_sanity.py}: the untrained model outputs exactly
zero). Attention uses PyTorch SDPA in both backbones; no fused-attention library is used.

\textbf{Training.} A single training loop drives both objectives. It supports torchrun DDP,
YAML configs, checkpoint resume, an exponential moving average of the weights (all reported
FID uses EMA weights), per-epoch EMA-only snapshots for quality-vs-progress curves, and cost
accounting: parameters, measured FLOPs, step time, throughput, peak memory and MFU are
written to \texttt{results/<run>.json} on every run.

\textbf{Samplers.} Ancestral DDPM sampling, deterministic DDIM sampling with a configurable
NFE, and Euler integration of the flow ODE. All three reuse the same trained \(\epsilon\)- or
velocity-predicting network, so they can be compared at fixed weights.

%=============================================================================
\section{Experimental setup}
\label{sec:setup}
%=============================================================================

\textbf{Hardware.} TODO: fill in the machine actually used (GPU model, interconnect,
NUMA layout). The topology matters: on the six-A6000 node used for the systems numbers,
GPUs 0--3 and 4--5 sit in different NUMA domains and the cross-domain link is an order of
magnitude slower, which is why \texttt{NCCL\_P2P\_DISABLE=1} is required and why DDP is kept
inside one domain.

\textbf{Protocol.} All cells of a scale share one frozen protocol: identical optimizer,
learning-rate schedule, batch size (global), EMA decay and seed (Table~\ref{tab:protocol}).
Unlike recipes that tune each cell separately, this design means a difference between cells
is attributable to the backbone or the objective rather than to the optimizer.

\begin{table}[t]
  \caption{The frozen protocol. Any deviation is recorded in the run's JSON.}
  \label{tab:protocol}
  \centering
  \small
  \begin{tabular}{@{}lll@{}}
    \toprule
    & MNIST $28^2$ & CelebA $64^2$ \\
    \midrule
    Cells & \multicolumn{2}{l}{UNet/DiT $\times$ $\epsilon$/flow, one run each} \\
    Epochs & 200 & 300 \\
    Global batch & 256 & 800 (160 per GPU $\times$ 5) \\
    Optimizer & \multicolumn{2}{l}{Adam, lr $2\times10^{-4}$, warmup then cosine to $10^{-6}$} \\
    Gradient clipping & \multicolumn{2}{l}{1.0} \\
    EMA decay & \multicolumn{2}{l}{0.9999, applied in every cell} \\
    Seed & \multicolumn{2}{l}{42} \\
    \bottomrule
  \end{tabular}
\end{table}

\textbf{FID protocol.} Reference statistics are computed from the training tensor itself
(\texttt{data/celeba\_64\_uint8.pt} for CelebA, the MNIST training split for MNIST) with the
same preprocessing chain used to evaluate samples; third-party statistics are never used,
because any difference in resizing makes the numbers incomparable. CelebA uses the standard
Inception-v3 pool3 features at $299^2$. Inception features are not meaningful for
grayscale digits, so the MNIST cells are evaluated in the feature space of a small
convolutional classifier trained on MNIST inside this repository
(Table~\ref{tab:protocol} footnote); that metric is \emph{non-standard} and only comparable
within this report. Tier~1 evaluations use 10k samples, Tier~2 (final) 50k, and the spread
is reported as the variability across five random 10k subsets.

%=============================================================================
\section{Sampler study (RQ2)}
\label{sec:samplers}
%=============================================================================

\begin{figure}[t]
  \centering
  \fig[width=\linewidth]{figures/fid_vs_nfe.pdf}
  \caption{Quality against the number of function evaluations, by dataset (note the
  logarithmic axes and the separate panels: the CelebA FIDs sit an order of magnitude above
  the MNIST ones). Within a panel every curve is a trained model plus a sampler; the DDPM
  vs.\ DDIM pair is a same-weights sampler ablation, while comparisons across cells also
  change the training objective.}
  \label{fig:fid_nfe}
\end{figure}

\subsection{What a better sampler buys, at CelebA scale}

The practical promise of DDIM is that a fixed trained model can be sampled in far fewer
network evaluations. Table~\ref{tab:celeba-samplers} answers that for the existing CelebA
DDPM checkpoint (UNet, 174.1M parameters, 300 epochs, trained before this protocol freeze).
Quality improves monotonically with the budget, and the 50-step deterministic sampler lands
within 12\% of the 1000-step ancestral sampler at $5.2\times$ the throughput: the last 950
evaluations of the ancestral chain buy 1.4 FID points and cost 5.9 hours per 50k-sample
evaluation. The 10-step sampler is the other end of the trade --- $100\times$ fewer
evaluations for $2.8\times$ the FID.

\begin{table}[t]
  \caption{CelebA $64^2$: one trained model, four sampling budgets (10k samples each, EMA
  weights). Throughput is measured on the same runs; the last column projects the cost of a
  50k-sample evaluation.}
  \label{tab:celeba-samplers}
  \centering
  \small
  \begin{tabular}{@{}lrrrr@{}}
    \toprule
    Sampler & NFE & FID & samples/s & 50k evaluation \\
    \midrule
    DDIM      & 10   & 31.40 & 49.5 & 17 min \\
    DDIM      & 20   & 19.13 & 24.8 & 34 min \\
    DDIM      & 50   & 12.75 &  9.9 & 1.4 h \\
    ancestral & 1000 & 11.37 &  1.9 & 7.3 h \\
    \bottomrule
  \end{tabular}
\end{table}

\begin{figure}[t]
  \centering
  \fig[width=2.65in]{figures/throughput_vs_nfe.pdf}
  \caption{Sampling throughput against NFE, measured on the same runs. Throughput is
  inversely proportional to NFE (log--log slope $\approx -1$ for every model), so the
  quality--NFE curves above are also quality--wall-clock curves up to a per-model constant.}
  \label{fig:throughput}
\end{figure}

\subsection{Quality--NFE frontier on MNIST}

Table~\ref{tab:mnist-nfe} is the same measurement for all four MNIST cells, which share a
training budget and protocol. Two patterns stand out. First, the rectified-flow cells
dominate $\epsilon$-prediction once the budget exceeds $\sim$16 evaluations: at 64
evaluations the flow cells reach FID $0.88$ (UNet) and $4.99$ (DiT), where the best
$\epsilon$-prediction result is $4.32$ at 50. The crossover sits near NFE 16, where the two
families are within 10\% of each other. Second, every curve is monotone in the budget ---
with one exception, the transformer's $\epsilon$ cell, which is the subject of the next
subsection.

\begin{table}[t]
  \caption{MNIST FID (10k samples, EMA weights) as a function of the sampling budget.
  ``---'' marks budgets that were not run for that cell. The 1000 column is the ancestral
  DDPM sampler; the other columns are DDIM (for $\epsilon$ cells) or Euler (for flow cells).}
  \label{tab:mnist-nfe}
  \centering
  \small
  \begin{tabular}{@{}llrrrrrrr@{}}
    \toprule
    Cell & Sampler & 10 & 16 & 20 & 32 & 50 & 64 & 1000 \\
    \midrule
    A: UNet $+$ $\epsilon$ & DDIM  & 22.06 & --- & 4.57 & --- & 4.32 & --- & 2.79 \\
    B: DiT $+$ $\epsilon$  & DDIM  & 49.14 & --- & 159.68 & --- & 672.39 & --- & 33.18 \\
    C: UNet $+$ flow       & Euler & --- & 5.06 & --- & 1.79 & --- & 0.88 & --- \\
    D: DiT $+$ flow        & Euler & --- & 12.98 & --- & 7.03 & --- & 4.99 & --- \\
    \bottomrule
  \end{tabular}
\end{table}

\subsection{The deterministic sampler collapses on the transformer}

Cell B is the one place in this study where more computation makes the result worse: its FID
\emph{grows} from 49.1 at NFE 10 to 159.7 at 20 and 672.4 at 50, while 1000-step ancestral
sampling of the same weights gives 33.2. Four checks localize the cause.

\textbf{Not the weights.} Sampling the raw (non-EMA) weights reproduces the failure
(FID 273.5 and 274.6 at NFE 20 and 50), so the exponential moving average is not at fault.

\textbf{Not the model.} The identical checkpoint with the stochastic ancestral sampler
reaches 33.2, and the training-time sample grids (ancestral, EMA weights) are recognizable
digits. The transformer learned the distribution well enough to be sampled.

\textbf{Not the sampler code.} The same function improves the UNet monotonically
($22.06 \rightarrow 4.57 \rightarrow 4.32$); its arithmetic does not depend on the backbone.

\textbf{It is the missing noise injection.} DDIM's $\eta$ interpolates between the
deterministic update ($\eta = 0$) and the ancestral posterior variance ($\eta = 1$). At fixed
NFE $=50$, raising $\eta$ restores the data statistics monotonically
(Table~\ref{tab:eta}) and turns the samples from saturated blobs back into digits.

\begin{table}[t]
  \caption{DiT $+$ $\epsilon$, DDIM at NFE $=50$, as a function of the noise-injection level
  $\eta$. ``Saturated'' is the fraction of pixels pinned at $\pm1$ after the final clamp;
  MNIST digits are near-binary, so the ancestral sampler (NFE 1000, $\eta=1$) saturates 84.8\%
  of pixels. Nothing else about the sampler changes.}
  \label{tab:eta}
  \centering
  \small
  \begin{tabular}{@{}lrrrrr@{}}
    \toprule
    $\eta$ & 0 & 0.25 & 0.50 & 0.75 & 1.00 \\
    \midrule
    saturated pixels & 47.7\% & 57.7\% & 66.0\% & 74.3\% & 79.1\% \\
    mean $|x|$       & 0.904  & 0.908 & 0.916 & 0.927 & 0.934 \\
    \bottomrule
  \end{tabular}
\end{table}

\textbf{Why this happens.} The deterministic recursion integrates the model's
$\epsilon$ error coherently along the trajectory, whereas each ancestral step re-injects the
schedule's noise and thereby bounds the drift. The transformer's error is large enough for
that difference to matter: measured against the same noised batches
(\texttt{eval/eps\_error\_by\_t.py}), its $\epsilon$ MSE is 30--55\% higher than the UNet's
across the whole trajectory, with the gap widest at low noise levels --- precisely where the
deterministic update reads its final image directly out of the model's $x_0$ estimate. The
same sampler that is the cheapest way to sample the UNet is the worst way to sample the
transformer, and the transformer is the model whose per-sample cost is lowest.

The practical reading is not that DDIM is fragile in general --- it is the best sampler in
Table~\ref{tab:celeba-samplers} and on the UNet --- but that the NFE reduction it offers is
conditional on score accuracy, and a controlled study is the only way to notice when the
condition fails. A paper reporting only the UNet numbers would have shipped a sampler
recommendation that inverts on the other backbone.

%=============================================================================
\section{Backbone and objective (RQ1)}
\label{sec:backbone}
%=============================================================================

\begin{table}[t]
  \caption{The 2$\times$2 design. FID at each cell's best NFE; see the text for how the
  NFE was chosen.}
  \label{tab:2x2}
  \centering
  \small
  \begin{tabular}{@{}llrr@{}}
    \toprule
    Dataset & Cell & FID & NFE at that FID \\
    \midrule
    MNIST & A: UNet + $\epsilon$   & \texttt{--} & \texttt{--} \\
    MNIST & B: DiT + $\epsilon$    & \texttt{--} & \texttt{--} \\
    MNIST & C: UNet + flow         & \texttt{--} & \texttt{--} \\
    MNIST & D: DiT + flow          & \texttt{--} & \texttt{--} \\
    CelebA & A: UNet + $\epsilon$  & \texttt{--} & \texttt{--} \\
    CelebA & B: DiT + $\epsilon$   & \texttt{--} & \texttt{--} \\
    CelebA & C: UNet + flow        & \texttt{--} & \texttt{--} \\
    CelebA & D: DiT + flow         & \texttt{--} & \texttt{--} \\
    \bottomrule
  \end{tabular}
\end{table}

\begin{figure}[t]
  \centering
  \fig[width=\linewidth]{figures/samples.png}
  \caption{EMA samples from each cell's final checkpoint.}
  \label{fig:samples}
\end{figure}

TODO: prose. State explicitly that parameter count is not compute: report the FLOPs column
from Table~\ref{tab:models} alongside every quality comparison, and note that the CelebA
cells are the expensive ones (measured training cost in Table~\ref{tab:cost}).

%=============================================================================
\section{Systems (RQ3)}
\label{sec:systems}
%=============================================================================

\begin{table}[t]
  \caption{Training step cost under each execution setting (30 steps, 5 warmups discarded).
  Produced by \texttt{eval/systems.py}; the FID column is the check that a speedup did not
  silently cost quality.}
  \label{tab:cost}
  \centering
  \small
  \begin{tabular}{@{}llrrrr@{}}
    \toprule
    Config & Setting & ms/step & img/s & peak GB & MFU \\
    \midrule
    TODO & FP32 & & & & \\
    TODO & BF16 & & & & \\
    TODO & compile & & & & \\
    TODO & BF16 + compile & & & & \\
    \bottomrule
  \end{tabular}
\end{table}

\begin{figure}[t]
  \centering
  \fig[width=2.65in]{figures/systems_settings.pdf}
  \caption{Step time by execution setting.}
  \label{fig:settings}
\end{figure}

\begin{figure}[t]
  \centering
  \fig[width=2.65in]{figures/profile_breakdown.pdf}
  \caption{Where GPU time goes in a diffusion training step, by kernel category. The
  measurement is a matmul/conv accounting only: elementwise-heavy operations
  (GroupNorm, SiLU, adaLN modulation) are exactly the ones a per-kernel profile reveals and
  a FLOPs count hides.}
  \label{fig:profile}
\end{figure}

TODO: the profile-first argument --- which category dominated, what optimization it implied,
and what the optimization actually returned. Then the distributed measurement: DDP step time
and the NCCL share, and the FSDP reverse benchmark at this parameter scale (expected to be
slower; the point is to measure and report it). Note the MFU caveat: the counter covers
matmul and convolution work only, so "MFU" here is a matmul-utilisation proxy.

%=============================================================================
\section{What didn't work}
\label{sec:failures}
%=============================================================================

Each entry follows the same shape: hypothesis, setup, observation, diagnosis, resolution,
lesson.

\textbf{Learning rate $5\times10^{-4}$ collapses CelebA training.} TODO: fill in with the
observed loss/grad-norm curves; resolution was $2\times10^{-4}$ with gradient clipping at
1.0.

\textbf{DataParallel deadlocks with a multi-worker DataLoader.} TODO: describe the hang and
the migration to torchrun-native DDP.

\textbf{P2P hangs on this node's NUMA layout.} TODO: NCCL\_P2P\_DISABLE=1, and what the
topology measurement showed.

\textbf{FSDP at this scale.} TODO: reverse benchmark numbers (expected: communication
overhead without meaningful memory savings at $\sim$130--174M parameters).

%=============================================================================
\section{Discussion and limitations}
\label{sec:limitations}
%=============================================================================

\begin{itemize}
  \item Unconditional generation only: no classifier-free guidance, no text conditioning.
  \item Two datasets, one resolution each; the CelebA $2\times2$ is the expensive half and
        its status is stated explicitly in each table.
  \item FID is computed against reference statistics built from this repository's own
        preprocessing chain. Absolute values are therefore \emph{not} comparable with
        published numbers; only comparisons within this report are meaningful.
  \item The MNIST metric is a non-standard feature-space distance.
  \item One seed per cell; the reported spread is FID variability across data subsets, not
        across seeds.
  \item FLOPs counts exclude elementwise work, so MFU figures are a proxy.
\end{itemize}

%=============================================================================
\section*{Reproduction}
%=============================================================================

All commands (environment, data, smoke tests, launching cells, evaluation) are in
\texttt{reports/runbook.md}; the raw archive behind every number is \texttt{results/*.json};
the committed code contains the full training stack and the test suite used to validate the
implementations (\texttt{tests/test\_sanity.py}, 14 CPU-only invariants).

\begin{thebibliography}{9}

\bibitem{ho2020ddpm}
Jonathan Ho, Ajay Jain, and Pieter Abbeel.
\newblock Denoising diffusion probabilistic models.
\newblock \emph{NeurIPS}, 2020.

\bibitem{song2020ddim}
Jiaming Song, Chenlin Meng, and Stefano Ermon.
\newblock Denoising diffusion implicit models.
\newblock \emph{ICLR}, 2021.

\bibitem{song2021sde}
Yang Song, Jascha Sohl-Dickstein, Diederik P. Kingma, Abhishek Kumar, Stefano Ermon, and
Ben Poole.
\newblock Score-based generative modeling through stochastic differential equations.
\newblock \emph{ICLR}, 2021.

\bibitem{nichol2021improved}
Alexander Quinn Nichol and Prafulla Dhariwal.
\newblock Improved denoising diffusion probabilistic models.
\newblock \emph{ICML}, 2021.

\bibitem{ho2022cfg}
Jonathan Ho and Tim Salimans.
\newblock Classifier-free diffusion guidance.
\newblock \emph{arXiv:2207.12598}, 2022.

\bibitem{rombach2022ldm}
Robin Rombach, Andreas Blattmann, Dominik Lorenz, Patrick Esser, and Bj\"orn Ommer.
\newblock High-resolution image synthesis with latent diffusion models.
\newblock \emph{CVPR}, 2022.

\bibitem{peebles2023dit}
William Peebles and Saining Xie.
\newblock Scalable diffusion models with transformers.
\newblock \emph{ICCV}, 2023.

\bibitem{liu2022flow}
Xingchao Liu, Chengyue Gong, and Qiang Liu.
\newblock Flow straight and fast: learning to generate and transfer data with rectified
flow.
\newblock \emph{ICLR}, 2023.

\bibitem{lipman2022fm}
Yaron Lipman, Ricky T. Q. Chen, Heli Ben-Hamu, Maximilian Nickel, and Matthew Le.
\newblock Flow matching for generative modeling.
\newblock \emph{ICLR}, 2023.

\bibitem{esser2024sd3}
Patrick Esser et al.
\newblock Scaling rectified flow transformers for high-resolution image synthesis.
\newblock \emph{ICML}, 2024.

\bibitem{heusel2017fid}
Martin Heusel, Hubert Ramsauer, Thomas Unterthiner, Bernhard Nessler, and Sepp Hochreiter.
\newblock GANs trained by a two time-scale update rule converge to a local Nash equilibrium.
\newblock \emph{NeurIPS}, 2017.

\bibitem{dao2022flash}
Tri Dao, Daniel Y. Fu, Stefano Ermon, Atri Rudra, and Christopher R\'e.
\newblock FlashAttention: fast and memory-efficient exact attention with IO-awareness.
\newblock \emph{NeurIPS}, 2022.

\bibitem{rajbhandari2020zero}
Samyam Rajbhandari, Jeff Rasley, Olatunji Ruwase, and Yuxiong He.
\newblock ZeRO: memory optimizations toward training trillion parameter models.
\newblock \emph{SC20}, 2020.

\bibitem{zhao2023fsdp}
Yanli Zhao, Andrew Gu, Rohan Varma, Liang Luo, Chien-Chin Huang, Min Xu, Less Wright,
Hamid Shojanazeri, Myle Ott, Sam Shleifer, et al.
\newblock PyTorch FSDP: experiences on scaling fully sharded data parallel.
\newblock \emph{VLDB}, 2023.

\end{thebibliography}

\end{document}
